% Exercise 3.1-2

\begin{proof}
  Since $b$ is an arbitrary positive real, the binomial theorem gives no finite
  expansion of ${(n + a)}^b$, and we instead bound $n + a$ above and below by
  constant multiples of $n$.

  Take $n_0 = 2\lvert a \rvert$, and let $n \geq n_0$. Then $\lvert a \rvert
  \leq n/2$, so
  \[
    \frac{n}{2} = n - \frac{n}{2} \leq n - \lvert a \rvert \leq n + a
    \leq n + \lvert a \rvert \leq n + \frac{n}{2} \leq 2n.
  \]
  Every quantity above is positive, and $t \mapsto t^b$ is monotonically
  increasing for $b > 0$, so raising the inequality to the power $b$ preserves
  it:
  \[
    {\biggl(\frac{n}{2}\biggr)}^b \leq {(n + a)}^b \leq {(2n)}^b,
  \]
  that is,
  \[
    0 \leq \frac{1}{2^b} n^b \leq {(n + a)}^b \leq 2^b n^b
  \]
  for all $n \geq n_0$.

  Thus with the identifications $c_1 = 2^{-b}$, $c_2 = 2^b$ and
  $n_0 = 2\lvert a \rvert$, we have ${(n + a)}^b = \Theta(n^b)$.
\end{proof}
